%% SYNTHETIC TEST UNIT generated by writing/tools/make_outline_book.py from FINAL_OUTLINE_v2 -- not book text
\chapter{阈值重求和与重新因子化}\label{ch:thr}
本单元为测试用的合成文本（由 FINAL\_OUTLINE\_v2 自动生成），陈述内容从略。

\section{第 1 节（测试）}\label{sec:thr-limit}

\begin{definition}[部分子阈值与强子阈值]\label{def:thr-limit}
（测试用合成陈述：内容从略。）
\uses{def:dy-kin, def:prelim-mellin}
\end{definition}

\begin{theorem}[DY 阈值重新因子化与指数化 \tagA]\label{thm:thr-refactorization}
（测试用合成陈述：内容从略。）
\uses{thm:dy-factorization, def:sudakov-cusp, prop:qcd-rge, def:thr-limit}
\end{theorem}
\begin{proofsketch}测试。\end{proofsketch}
\paragraph{地位与缺口}测试：本陈述的地位由\cref{thm:dy-factorization}、\cref{prop:disall-two-routes}决定。

\section{第 2 节（测试）}\label{sec:thr-refactorization}

\begin{proposition}[单圈检验 \tagR]\label{prop:thr-oneloop}
（$R_{\mathrm{PT}}$：在重整化微扰论中对红外正规化的部分子态严格。）
（测试用合成陈述：内容从略。）
\uses{prop:dy-nlo, lem:prelim-mellin, lem:appA-mellin-table, def:thr-limit}
\end{proposition}

\begin{remark}[Landau 极点与处方]\label{rem:thr-prescriptions}
（测试用合成陈述：内容从略。）
% uses: none
\end{remark}

\section{第 3 节（测试）}\label{sec:thr-check}

\begin{proposition}[DIS 的大 $x$ 重新因子化 \tagA]\label{prop:thr-dis}
（测试用合成陈述：内容从略。）
\uses{thm:disall-factorization, def:thr-limit}
\end{proposition}
\begin{proofsketch}测试。\end{proofsketch}
\paragraph{地位与缺口}测试：本陈述的地位由\cref{thm:disall-factorization}、\cref{prop:disall-two-routes}决定。

\begin{proposition}[次领头幂次的阈值对数 \tagP]\label{prop:thr-nlp}
（测试用合成陈述：内容从略。）
\uses{thm:thr-refactorization}
\end{proposition}
\paragraph{地位与缺口}测试：本陈述的地位由\cref{prop:power-nlp-endpoint}决定。

\section{第 4 节（测试）}\label{sec:thr-prescriptions}

\begin{example}[指数的 $O(\alpha_s)$ 展开与 prop:dy-nlo 的 Mellin 变换]\label{ex:thr-dy}
（测试用合成陈述：内容从略。）
\uses{prop:thr-oneloop}
\end{example}
